\chapter{Aprendizaje Continuo}
\label{cap:aprendizaje}

Los parámetros del modelo (los tiempos por tramo, las proporciones por longitud y los factores de ritmo) no se
escriben a mano: se aprenden del propio feed en tiempo real. Este capítulo describe el ciclo completo: cómo se captura
el feed, cómo se reconstruye el instante real en que cada autobús pasa por cada parada, qué muestras se extraen, cuánto
se conservan y con qué reglas se decide si una nueva calibración sustituye a la anterior.

\section{El Ciclo}
\label{sec:ciclo}

\begin{enumerate}
  \item \textbf{Captura}, continua y en bloques de una hora de reloj: se descarga el feed cada 20~s y se guarda cada
        versión distinta.
  \item \textbf{Muestras y medida}, al terminar cada bloque: se reconstruyen los pasos por parada, se extraen las
        muestras de tramo y de ritmo y se añaden a dos ficheros acumulados; además se mide la precisión de las
        llegadas (capítulo «\nameref{cap:evaluacion}») y se archivan las capturas en bruto del bloque.
  \item \textbf{Recalibración}, cada día a las 04:00 UTC: con todo lo acumulado se calculan nuevos parámetros, se
        validan y, solo si mejoran y cambian, se escriben.
  \item \textbf{Despliegue}, cada día a las 03:00 UTC: la aplicación se reconstruye con la última versión de los
        parámetros. Producción usa, por tanto, parámetros de entre 23 y 47~horas de antigüedad, salvo despliegue manual.
\end{enumerate}

Todo el ciclo usa el mismo código de localización y predicción que producción (la misma clase), de modo que lo que se
aprende y lo que se mide es lo que ve el usuario. La captura se ejecuta en una plataforma de integración continua que
limita cada trabajo a seis horas y que retrasa o se salta muchos trabajos programados: hasta el 31 de octubre de 2026
la captura era de 45~minutos cada hora y solo se ejecutaron 163 de unas 600 capturas previstas, algunas con horas de
retraso. Desde entonces cada trabajo captura unas cinco horas y media y, si termina bien, lanza el siguiente, que
espera su turno y empieza en cuanto acaba el anterior; un disparo programado cada hora queda como respaldo.

\section{Captura del Feed}
\label{sec:captura}

Cuando se fijó el intervalo de descarga, en octubre de 2026, el servidor del feed publicaba una versión nueva cada
55~s aproximadamente (el 5 de noviembre, cada 64~s de mediana). Cada 20~s (\texttt{intervalo\_s}) se
descarga, y la versión se guarda solo si su marca \texttt{RecordedAtTime} es distinta de la anterior. Con el intervalo
inicial de 45~s se perdía alguna versión cuando dos publicaciones caían entre dos descargas (46 versiones guardadas en
45~minutos, de unas 49 publicadas); con 20~s se guardan todas, y muestrear más deprisa no aporta nada, porque el feed
no cambia entre publicaciones.

La captura se organiza en bloques que terminan a la hora en punto, de modo que ninguno cruza la medianoche y cada uno
usa el horario de su día. Las capturas no usan el emparejador de ayer de la sección «\nameref{sec:mas-de-24h}»: la
parte de un viaje posterior a la medianoche no se sitúa con el GPS, y no aporta muestras ni entra en las medidas. En
Bizkaibus afecta sobre todo a los 11 viajes de fin de semana que acaban después de las 29:00, porque el resto de la
noche el feed no se actualiza. Mientras se captura un bloque se procesa el anterior, sin huecos entre ellos. Para medir la
precisión de las predicciones de un bloque se usan también las capturas del bloque siguiente, solo para ver cuándo
pasan los autobuses que estaban a punto de llegar al terminar el bloque. Las capturas en bruto de cada bloque,
comprimidas (unos 250~KB por hora), se archivan junto con las muestras nuevas que salieron de él, sin límite de
antigüedad: cualquier cálculo se puede repetir más adelante aunque los ficheros acumulados descarten lo antiguo.

Al procesar una captura, cada fichero se identifica por su \texttt{RecordedAtTime}
(los duplicados se ignoran) y se convierte a fecha y segundos de Madrid. Las instantáneas se ordenan por tiempo y se
agrupan en \emph{ventanas de observación}: una ventana termina cuando el hueco hasta la siguiente instantánea supera
300~s (\texttt{MAX\_GAP\_BETWEEN\_SNAPSHOTS\_SECONDS}). Cada instantánea guarda el instante final de su ventana,
$t_{\mathrm{fin}}$.

\section{Pistas}
\label{sec:pistas}

Para cada instantánea, en orden, se construye el emparejador con $\tnow$ igual a la hora de la instantánea y la fecha
de la instantánea como fecha de servicio, y para cada entrada del feed con viaje conocido se calcula su posición
(capa 1, con el contraste de la sección «\nameref{sec:desconfianza}»). Si existe, se añade el punto $(\tgps,\Tpos)$ a la
\emph{pista} de la clave (vehículo, viaje).

\begin{definicion}[Pista]
La pista de un vehículo $v$ en un viaje $J$ durante una captura es la sucesión, en el orden de las instantáneas, de
los puntos $(t_k,\theta_k)$, con $t_k=\tgps$ y $\theta_k=\Tpos$ de cada instantánea en que el vehículo tenía posición
en ese viaje.
\end{definicion}

Solo se aprende de posiciones GPS válidas: las capas 2 y 3, que son estimaciones, nunca alimentan el aprendizaje, y
tampoco las posiciones en la cabecera (sección «\nameref{sec:cabecera}»), que no describen ningún recorrido.

\section{El Instante Real de Paso por una Parada}
\label{sec:paso-real}

No existe en el feed ningún dato de «el autobús pasó por esta parada». Se reconstruye a partir de la pista, usando el
horario como escala de distancia: cada $\theta_k$ indica, en horas del horario, dónde estaba el autobús.

\begin{definicion}[Paso real]
\label{def:paso-real}
Dada una pista $\big((t_k,\theta_k)\big)_{k=1}^{K}$ y la hora publicada $a$ de una parada, el paso real $P(a)$ se
obtiene recorriendo los pares consecutivos $k=2,\dots,K$ y tomando el primero que cumpla
\[
\theta_k>\theta_{k-1},\qquad t_k-t_{k-1}\le300\ \text{s},\qquad \theta_{k-1}\le a\le\theta_k;
\]
entonces
\begin{equation}
\label{eq:paso-real}
P(a)=t_{k-1}+\frac{a-\theta_{k-1}}{\theta_k-\theta_{k-1}}\,\big(t_k-t_{k-1}\big).
\end{equation}
Si ningún par cumple las condiciones, el paso real no se conoce.
\end{definicion}

Es una interpolación lineal entre las dos observaciones que rodean a la parada. Los pares en que el autobús no avanza
($\theta_k\le\theta_{k-1}$, por ejemplo una posición repetida) o separados más de cinco minutos se descartan, porque no
permiten interpolar con fiabilidad.

\begin{ejemplo}
Una pista $(28\,900,\,28\,880)$, $(28\,960,\,28\,940)$, $(29\,020,\,29\,030)$ y una parada con $a=28\,950$: el par que
la rodea es el segundo, y $P=28\,960+\frac{28\,950-28\,940}{29\,030-28\,940}\cdot60=28\,966{,}67$~s. Para una parada con
$a=28\,800$, anterior a la primera observación, el paso no se conoce.
\end{ejemplo}

Con una posición por minuto, el error de la interpolación está acotado por el tiempo entre dos
observaciones: si el autobús se detiene o acelera entre ellas, el paso real puede estar en cualquier punto del
intervalo. En la práctica se trata como una resolución de $\pm30$~s, que es también la del método independiente que se
usa para evaluar (sección «\nameref{sec:verdad}»).

\section{Muestras de Tramo}
\label{sec:muestras-tramo}

Para cada pista y cada par de paradas consecutivas $(p_i,p_{i+1})$ de su viaje, si se conocen $P_i=P(a_i)$ y
$P_{i+1}=P(a_{i+1})$ y $P_{i+1}>P_i$, se produce una muestra de tramo con los campos
\begin{center}
\small
\begin{tabular}{@{}ll@{}}
\toprule
Campo & Valor\\
\midrule
\texttt{track} & «$D|v|J$»: fecha, vehículo y viaje\\
\texttt{line}, \texttt{date} & línea y fecha de servicio\\
\texttt{at} & $\lfloor P_i\rfloor$ (truncado: parte entera)\\
\texttt{band} & $\franja(D,\lfloor P_i\rfloor)$\\
\texttt{from}, \texttt{to} & identificadores de $p_i$ y $p_{i+1}$\\
\texttt{scheduled} & $a_{i+1}-a_i$\\
\texttt{actual} & $\redondeo_1(P_{i+1}-P_i)$, con un decimal\\
\bottomrule
\end{tabular}
\end{center}

\noindent La franja es la del momento en que el autobús entra en el tramo. La pista incluye la fecha para que el mismo
vehículo en el mismo viaje en dos días distintos cuente como dos pistas.

\section{Muestras de Ritmo}
\label{sec:muestras-ritmo}

Las muestras de ritmo sirven para estimar $k$ (sección «\nameref{sec:k}»). En cada instantánea, para cada entrada con
posición y cada parada de su viaje con $a_\dest>\Tpos$, se registra la predicción: la hora de la posición, $\Tpos$ y
$a_\dest$. Cuando se conocen las pistas completas, se calcula el paso real $P(a_\dest)$ y la predicción se convierte en
muestra si cumple:
\begin{enumerate}
  \item el paso real se conoce;
  \item la ventana de observación se extiende al menos 10~minutos después de la hora publicada:
        $t_{\mathrm{fin}}\ge a_\dest+600$ (\texttt{OBSERVED\_AFTER\_PREDICTION\_SECONDS});
  \item el tiempo programado que faltaba es de al menos dos minutos: $H=a_\dest-\Tpos\ge120$~s
        (\texttt{MIN\_REMAINING\_SECONDS\_FOR\_FACTOR});
  \item el tiempo real que faltaba es positivo: $Y=P(a_\dest)-\tgps>0$.
\end{enumerate}
Se guardan la pista, el vehículo, la línea, la fecha, $\texttt{at}=\tgps$, la franja de $\tgps$, $H$, $Y$ y el paso
real (estos dos con un decimal).

\subsection{Por Qué Exigir la Ventana} Solo se conoce el paso real de los autobuses que pasan antes de que termine la
captura. Si se aceptaran todas las predicciones con paso conocido, las de autobuses lentos (que pasan después del
final) quedarían fuera y $k$ saldría sesgado hacia autobuses rápidos: es un problema de \emph{censura}. Exigir que la
ventana cubra al menos diez minutos después de la hora publicada reduce ese sesgo a los autobuses que van más de diez
minutos tarde respecto a su horario en ese tramo. La misma regla, con la llegada estimada en lugar de la publicada,
se aplica al medir el error (capítulo «\nameref{cap:evaluacion}»).

\subsection{Aclarado} Un autobús visto durante 45~minutos con veinte paradas por delante produce cientos de
predicciones muy correlacionadas. Para que no pese más que otros, se conservan como mucho 30 por pista
(\texttt{MAX\_SAMPLES\_PER\_TRACK}): si una pista tiene $c>30$ muestras, en su orden original, se conservan las de
índices
\begin{equation}
\label{eq:aclarado}
\Big\lfloor\frac{i\,c}{30}\Big\rfloor,\qquad i=0,1,\dots,29 .
\end{equation}
Con $c=100$ son los índices 0, 3, 6, 10, 13, 16, 20, \dots, 96: un muestreo uniforme que conserva toda la duración de
la pista.

\section{Acumulación y Retención}
\label{sec:retencion}

Las muestras se añaden a dos ficheros JSON Lines acumulados (formatos en el Apéndice~\ref{ap:formatos}). Al añadir:

\begin{itemize}
  \item se descartan las muestras con fecha anterior a 120~días antes de hoy (\texttt{SAMPLES\_RETENTION\_DAYS});
  \item en las muestras de tramo, si un tramo tiene más de 300 de un mismo tipo de día, laborable o fin de semana
        (\texttt{MAX\_SAMPLES\_PER\_SEGMENT}), se conservan las 300 más recientes de ese tipo por (fecha, \texttt{at}),
        sin alterar el orden del resto del fichero;
  \item en las de ritmo, si el total pasa de $500\,000$ (\texttt{MAX\_PACE\_SAMPLES}), se ordenan por (fecha,
        \texttt{at}) y se conservan las $500\,000$ más recientes.
\end{itemize}

\noindent Estos topes mantienen los ficheros por debajo de unos 120~MB y la recalibración por debajo de 2~GB de memoria,
con un margen amplio respecto a los límites de la plataforma. El tope hace además de memoria: al quedarse solo lo más
reciente, un cambio en la calle (un desvío, unas obras) entra en lo aprendido en uno o dos días.

Hasta el 5 de noviembre de 2026 el tope era de 300 muestras por tramo, sin separar el tipo de día. Con las primeras
capturas, 300 muestras eran unos diez días de un tramo muy transitado, y diez días siempre incluyen un fin de semana: se
comprobó con 25~días que en los 109 tramos (de 2\,487) que llegaban al tope las franjas estaban igual de repartidas que
en el resto. Desde el 2 de noviembre las capturas recogen unas diez veces más pasos al día, y el 5 de noviembre 533 de
los 2\,757 tramos estaban en el tope: muchos guardaban solo los últimos 3 a 5~días, y 127 no tenían ninguna muestra de
fin de semana. En esos tramos el sábado y el domingo se usaba el tiempo de laborable, y el lunes, al revés, la lista
estaba llena de muestras del fin de semana.

Con un tope por tipo de día cada lista se renueva por su cuenta y siempre queda un fin de semana guardado. Para medirlo
se aprendió con las muestras anteriores a cada día y se probó en ese día (Tabla~\ref{tab:tope-tipo-dia}). Con el tope
de 300 no se puede comparar con los datos guardados, porque el fichero ya venía recortado por el tope anterior; por eso
se reprodujo la misma situación con topes menores en los tramos que conservan todo su historial. El tope por tipo de día
nunca empeora y mejora sobre todo el lunes que sigue al fin de semana.

\begin{table}[ht]
\centering
\caption{Error por Tramo con el Tope por Tramo y con el Tope por Tramo y Tipo de Día}\label{tab:tope-tipo-dia}
\begin{tabular}{@{}llrrr@{}}
\toprule
Tope & Día de prueba & Pasos & Error medio (s) & Dentro de $\pm30$~s (\%)\\
\midrule
20 & sábado 31 oct. & 22\,787 & 19,3 $\to$ 19,0 & 82,8 $\to$ 83,3\\
20 & domingo 1 nov. & 21\,441 & 19,1 $\to$ 18,9 & 83,3 $\to$ 83,6\\
20 & lunes 2 nov. & 52\,910 & 21,4 $\to$ 20,3 & 79,4 $\to$ 81,1\\
20 & martes 3 nov. & 58\,509 & 20,9 $\to$ 20,9 & 80,4 $\to$ 80,4\\
40 & sábado 31 oct. & 22\,787 & 19,0 $\to$ 19,0 & 83,3 $\to$ 83,5\\
40 & domingo 1 nov. & 21\,543 & 18,9 $\to$ 18,7 & 83,7 $\to$ 84,0\\
40 & lunes 2 nov. & 52\,910 & 20,7 $\to$ 19,9 & 80,5 $\to$ 81,7\\
40 & martes 3 nov. & 61\,081 & 20,3 $\to$ 20,3 & 81,0 $\to$ 81,1\\
80 & sábado 31 oct. & 22\,787 & 19,0 $\to$ 19,0 & 83,5 $\to$ 83,5\\
80 & domingo 1 nov. & 21\,543 & 18,8 $\to$ 18,7 & 83,8 $\to$ 84,0\\
80 & lunes 2 nov. & 52\,910 & 20,0 $\to$ 19,8 & 81,7 $\to$ 82,0\\
80 & martes 3 nov. & 61\,081 & 20,1 $\to$ 20,0 & 81,5 $\to$ 81,6\\
\bottomrule
\end{tabular}

\nota{Error de cada tramo entre dos paradas seguidas, no de la llegada completa: antes de la flecha, con un tope por
tramo; después, con un tope por tramo y tipo de día. Se aprendió con las muestras anteriores al día de prueba, con las
mismas medianas por franja de la sección «\nameref{sec:contraccion-tramo}». Solo cuentan los pasos con tiempo aprendido en los dos
casos.}
\end{table}

Para que una red recién puesta en marcha no empiece sin datos, el repositorio del motor incluye una semilla de muestras
de tramo que se usa si todavía no existe el fichero acumulado.

\section{Recalibración Diaria}
\label{sec:recalibracion}

Cada día se recalibran, en este orden, el factor de ritmo, los tiempos por tramo y la corrección por viaje. El orden
importa: la validación de los tiempos por tramo compara contra el factor de ritmo que se acaba de escribir, y la
corrección por viaje se mide contra el horario corregido que dan los tiempos recién escritos.

\subsection{Pliegues}

La validación es cruzada en dos pliegues por \emph{pistas}: cada muestra pertenece al pliegue
\begin{equation}
\label{eq:pliegue}
\mathrm{pliegue}(x)=\mathrm{crc32}\big(\texttt{track}(x)\big)\bmod2,
\end{equation}
donde $\mathrm{crc32}$ es el CRC de 32 bits de IEEE 802.3 (el de \texttt{zlib}, PNG o Ethernet) aplicado a los bytes
UTF-8 de la cadena, interpretado \emph{sin signo}. Todas las muestras de una misma pista caen en el mismo pliegue, de
modo que nunca se aprende y se mide con el mismo recorrido (un autobús en un viaje de un día). El mismo vehículo, en
otro viaje u otro día, puede caer en el otro pliegue, y los dos pliegues comparten los días. Se aprende con un pliegue
y se mide en el otro, y al revés.

\subsection{Factor de Ritmo}

\begin{enumerate}
  \item Se ajustan los factores con todas las muestras (ecuación~\eqref{eq:k-contraido} y los mínimos de la
        sección «\nameref{sec:k}»).
  \item Para cada pliegue $f$, se ajustan los factores con las muestras del otro pliegue. Después, para cada muestra
        de $f$ con factor aplicable $k\ne1$, se calculan los errores absolutos en minutos
        \[
        e_{\mathrm{antes}}=\frac{|\texttt{at}+H-P|}{60},\qquad e_{\text{después}}=\frac{|\texttt{at}+k\,H-P|}{60}.
        \]
  \item Se escribe el fichero de factores solo si hay alguna muestra de prueba, si la media de $e_{\text{después}}$
        es estrictamente menor que la de $e_{\mathrm{antes}}$, y si los factores nuevos no son iguales a los del fichero
        actual.
\end{enumerate}

\subsection{Tiempos por Tramo y Proporciones}

\begin{enumerate}
  \item Para cada pliegue $f$: se ajustan los tiempos por tramo con el otro pliegue y, para cada muestra de $f$ cuyo
        tramo tenga tiempo aprendido en su franja (ecuación~\eqref{eq:consulta-tramo}), se calculan
        \[
        e_{\mathrm{antes}}=\frac{|k\,h-y|}{60},\qquad e_{\text{después}}=\frac{|\tau-y|}{60},
        \]
        con $k$ el factor de ritmo vigente para la línea, la fecha y el instante de la muestra.
  \item Del mismo modo se validan las proporciones por longitud (con la longitud mediana de cada par de paradas en el
        horario): para cada muestra de prueba con $h>0$ y respaldo definido, $|k\,h-y|$ frente a
        $|\tau^{\mathrm{L}}-y|$. Las proporciones se incluyen en la tabla solo si la suma de errores baja.
  \item Se escribe el fichero solo si la suma de $e_{\text{después}}$ de los tiempos por tramo es estrictamente
        menor que la de $e_{\mathrm{antes}}$ y si los tiempos o las proporciones cambian respecto al fichero actual. Si
        los tiempos por tramo no mejoran, no se escribe nada, aunque mejoren las proporciones.
\end{enumerate}

\noindent La línea de base de la validación es siempre lo que haría el sistema sin lo aprendido (el horario por el
factor de ritmo). Con las 12\,755 medidas iniciales, el error medio por tramo bajó de 0,85~min a 0,39~min en las
medidas de prueba cuyo tramo estaba aprendido; con las proporciones por longitud, los tramos sin aprender pasaron de
0,88 a 0,69~min. La Tabla~\ref{tab:curva} cuenta todas las medidas de prueba, aprendidas o no, y por eso parte de
0,89~min.

\subsection{Corrección por Viaje}

\begin{enumerate}
  \item Con las muestras de tramo se reconstruye el paso real $R$ de cada viaje por cada parada y día: una muestra da
        el paso por su parada de salida (\texttt{at}) y, si falta el de la de llegada, también ese (\texttt{at} más el
        tiempo real del tramo). La pista dice el viaje, que se busca en el horario o en sus alias. Por eso el
        recalibrado reconstruye antes la base del horario del día.
  \item Para cada paso se calcula el desvío $e=R-C_0$, donde $C_0$ es el horario corregido~\eqref{eq:capa3} sin
        $\Delta_v$, con la $D'$ de ese día. Se saltan la cabecera, las paradas repetidas en un recorrido y los pasos a
        más de 45~minutos de su hora publicada. Cada desvío se recorta a $\pm900$~s para que un día de atasco no mande.
  \item Con $b$ la franja de la salida de cabecera según $D'$, se calculan
        \begin{equation}
        \label{eq:correccion-tablas}
        \begin{gathered}
        g(b)=Q_{0{,}4}(e_b),\qquad
        l(\ell,b)=\frac{n_{\ell,b}\,Q_{0{,}4}(e_{\ell,b})+20\,g(b)}{n_{\ell,b}+20},\\
        \Delta(v,\theta,p)=\frac{\sum_{\text{días}}Q_{0{,}4}(e_{\text{día}})+l(\ell,b)}{n_{\text{días}}+1},
        \end{gathered}
        \end{equation}
        donde $Q_{0{,}4}$ es el percentil 40 (con la misma regla que la mediana baja: el valor de la posición
        $\lfloor0{,}4(n-1)\rfloor$ de los ordenados), $e_b$ los desvíos de la franja, $e_{\ell,b}$ los de la línea en
        ella, $n_{\ell,b}$ su número y $e_{\text{día}}$ los del viaje en ese tercio cada día. Un tercio sin pasos toma el
        valor de su línea. Todo se guarda en segundos enteros. Hasta el 3 de noviembre de 2026 eran medias; el 40 es el
        percentil de partida, y la recalibración elige cada noche el que más acierta (sección
        «\nameref{sec:correccion-percentil}»).
  \item Se aprende con todos los días menos los 7 últimos y se mide en esos 7 (una sola partición; la evaluación de la
        sección «\nameref{sec:eval-correccion}» es más exigente y aprende cada día de prueba con todos los anteriores),
        con cada percentil entre el 30 y el 50;
        se queda el que más acierta a $\pm2$~minutos. La tabla, aprendida con todos los días y ese percentil,
        se escribe solo si el error medio absoluto de la capa 3 baja con la corrección y si cambia respecto a la
        actual. La herramienta informa además de cada festivo de las muestras, aprendiendo con los demás días.
\end{enumerate}

\noindent Aquí la validación es por días y no por pistas: la corrección es de cada viaje, y con pliegues por pistas
el mismo viaje de otro día estaría en las dos mitades.

\subsection{Qué Garantiza y Qué No}

Las reglas de escritura son conservadoras: un día con datos anómalos (una avería del feed, un festivo atípico) solo
cambia el modelo si el cambio mejora la predicción en recorridos que no se usaron para aprenderlo. Si no hay cambios,
no se escribe nada y el fichero conserva su fecha. La validación por pistas es una \emph{salvaguarda} diaria, no un
método para elegir reglas: para decidir si una regla nueva entra en el método se usa la evaluación por días del
capítulo «\nameref{cap:evaluacion}», que reproduce mejor lo que pasará con datos futuros.

\section{Curva de Aprendizaje}
\label{sec:curva}

Para saber cuánto mejora el modelo con más datos, se separó la mitad de los autobuses de la primera sesión de
captura (12\,755 medidas de tramo, un laborable a mediodía) para probar, y con la otra mitad se aprendieron los tiempos
por tramo usando cada vez más autobuses (Tabla~\ref{tab:curva}).

\begin{table}[ht]
\centering
\caption{Curva de Aprendizaje con la Primera Sesión de Captura}\label{tab:curva}
\begin{tabular}{@{}rrrrr@{}}
\toprule
Autobuses & Tramos vistos & Cobertura & Aprendido y $k$ & Solo $k\cdot h$\\
\midrule
16  & 278   & 1,2\,\%  & 0,89 & 0,89\\
45  & 610   & 13,0\,\% & 0,84 & 0,89\\
92  & 1\,001 & 32,2\,\% & 0,77 & 0,89\\
196 & 1\,458 & 62,8\,\% & 0,62 & 0,89\\
379 & 1\,916 & 77,5\,\% & 0,52 & 0,89\\
\bottomrule
\end{tabular}
\nota{Error medio por tramo, en minutos, en los autobuses de prueba. Tramos vistos: tramos distintos en las medidas
usadas para aprender (aprendidos o no). Cobertura: parte de las medidas de prueba cuyo tramo tiene tiempo aprendido en
su franja; las demás usan $k\cdot h$.}
\end{table}

\begin{figure}[ht]
\centering
\caption{Error por Tramo Según el Número de Autobuses}\label{fig:curva}
\begin{tikzpicture}
\begin{axis}[
  width=0.85\textwidth, height=5.5cm,
  xmode=log, log basis x=2,
  xlabel={Autobuses usados para aprender (escala logarítmica)},
  ylabel={Error medio por tramo (min)},
  xtick={16,45,92,196,379}, xticklabels={16,45,92,196,379},
  ymin=0.45, ymax=0.95,
  yticklabel style={/pgf/number format/fixed,/pgf/number format/precision=2,/pgf/number format/use comma},
  grid=major, grid style={gris!20},
  legend style={at={(0.02,0.05)},anchor=south west,font=\footnotesize,draw=none,fill=none},
  every axis plot/.append style={thick}]
\addplot[ruta,mark=*] coordinates {(16,0.89) (45,0.84) (92,0.77) (196,0.62) (379,0.52)};
\addlegendentry{Lo aprendido y, si no, $k\cdot h$}
\addplot[gris,dashed] coordinates {(16,0.89) (379,0.89)};
\addlegendentry{Solo $k\cdot h$}
\end{axis}
\end{tikzpicture}
\nota{Las medidas de prueba sin tramo aprendido (100\,\% menos la cobertura de la Tabla~\ref{tab:curva}) bajan del
98,8 al 22,5\,\%, un 40--45\,\% en cada una de las dos últimas duplicaciones, y el error sigue bajando.}
\end{figure}

El error no se había estabilizado con 379 autobuses. Desde entonces la captura pasó a ser horaria, y a finales de
octubre de 2026 las muestras acumuladas suman 195\,786 tramos de 18\,604 pistas en 26 días, con 2\,454 tramos aprendidos.
